% GENERATED FILE. Edit canonical lessons, then run npm run generate:books.
% canonical-source-hash: fnv1a64:6725545e15619bb9
% canonical-lessons: MR-C512-karkhana, MR-C512-utpadan, MR-C512-mal, MR-C512-puravtha, MR-C512-khapne

\chapter{Made in a Factory}
\label{ch:mr-made-in-a-factory}
\hlchaptermodality{512}

\begin{chapteropening}
\textbf{By the end of this chapter:} \emph{I can talk about factories, goods and supplies.}

Complete the last lesson of chapter 512: I can talk about factories, goods and supplies.
\end{chapteropening}

\section[\texorpdfstring{\mr{कारखाना} (kārkhānā)}{कारखाना (kārkhānā)}]{\mr{कारखाना} (kārkhānā) — a factory}

\label{lesson:MR-C512-karkhana}



\begin{quote}
Before the new one: say the Marathi for a business partner, then the Marathi for to invest.
\end{quote}

\subsection*{You'll want to know: \mr{कारखाना}}
\textbf{\mr{कारखाना}} — \emph{kārkhānā} — \textquotedblleft{}a factory\textquotedblright{}.

\begin{quote}
\mr{आमच्या} \mr{शहराजवळ} \mr{साखरेचा} \mr{कारखाना} \mr{आहे}.
\end{quote}

\begin{quote}
\emph{āmcyā śahrājvaḷ sākhrecā kārkhānā āhe.}
\end{quote}

\begin{quote}
There is a sugar factory near our town.
\end{quote}

Masculine. Plural: \mr{कारखाने}, \emph{kārkhāne}.

\begin{cousinweb}
\textbf{\mr{कारखाना}} (\emph{kārkhānā}) is a loan from Persian \emph{kārkhāna}, \textquotedblleft{}workshop\textquotedblright{}, from \emph{kār}, \textquotedblleft{}work\textquotedblright{}, and \emph{khāna}, \textquotedblleft{}house\textquotedblright{}.
\end{cousinweb}

\begin{grammarlens}[title={What you've built}]
One more word for describing what you lived and felt.
\end{grammarlens}

\subsection*{Your turn}
\begin{itemize}
\raggedright
  \item \emph{Say it:} \emph{kārkhānā}
  \item \emph{Say it:} \emph{kārkhānā}, once more
  \item \emph{Say it:} \emph{āmcyā śahrājvaḷ sākhrecā kārkhānā āhe.} — There is a sugar factory near our town.
  \item \emph{Recall:} say the Marathi for a business partner, then the Marathi for to invest, then say \emph{kārkhānā} again
\end{itemize}

\begin{tcolorbox}[breakable,colback=teal!4,colframe=teal!35!black,title={Before you move on}]
What does \mr{कारखाना} mean? (\textquotedblleft{}A factory\textquotedblright{}.) Say it once more.
\end{tcolorbox}
\section[\texorpdfstring{\mr{उत्पादन} (utpādan)}{उत्पादन (utpādan)}]{\mr{उत्पादन} (utpādan) — production; output}

\label{lesson:MR-C512-utpadan}



\begin{quote}
Before the new one: say the Marathi for to invest, then the Marathi for a factory.
\end{quote}

\subsection*{You'll want to know: \mr{उत्पादन}}
\textbf{\mr{उत्पादन}} — \emph{utpādan} — \textquotedblleft{}production; output\textquotedblright{}.

\begin{quote}
\mr{या} \mr{कारखान्यात} \mr{रोज} \mr{हजार} \mr{सायकलींचं} \mr{उत्पादन} \mr{होतं}.
\end{quote}

\begin{quote}
\emph{yā kārkhānyāt roj hajār sāyaklīñcaṁ utpādan hotaṁ.}
\end{quote}

\begin{quote}
A thousand bicycles are made in this factory every day.
\end{quote}

Neuter.

\mr{उत्पादन} \mr{होणे}, to be produced.

\begin{cousinweb}
\textbf{\mr{उत्पादन}} (\emph{utpādan}) is a tatsama, taken from Sanskrit \emph{utpādana}, \textquotedblleft{}bringing forth, producing\textquotedblright{} (Monier-Williams).
\end{cousinweb}

\begin{grammarlens}[title={What you've built}]
One more word for describing what you lived and felt.
\end{grammarlens}

\subsection*{Your turn}
\begin{itemize}
\raggedright
  \item \emph{Say it:} \emph{utpādan}
  \item \emph{Say it:} \emph{utpādan}, once more
  \item \emph{Say it:} \emph{yā kārkhānyāt roj hajār sāyaklīñcaṁ utpādan hotaṁ.} — A thousand bicycles are made in this factory every day.
  \item \emph{Recall:} say the Marathi for to invest, then the Marathi for a factory, then say \emph{utpādan} again
\end{itemize}

\begin{tcolorbox}[breakable,colback=teal!4,colframe=teal!35!black,title={Before you move on}]
What does \mr{उत्पादन} mean? (\textquotedblleft{}Production; output\textquotedblright{}.) Say it once more.
\end{tcolorbox}
\section[\texorpdfstring{\mr{माल} (māl)}{माल (māl)}]{\mr{माल} (māl) — goods; stock}

\label{lesson:MR-C512-mal}



\begin{quote}
Before the new one: say the Marathi for a factory, then the Marathi for production.
\end{quote}

\subsection*{You'll want to know: \mr{माल}}
\textbf{\mr{माल}} — \emph{māl} — \textquotedblleft{}goods; stock\textquotedblright{}.

\begin{quote}
\mr{कारखान्याचा} \mr{माल} \mr{रोज} \mr{बाजारात} \mr{जातो}.
\end{quote}

\begin{quote}
\emph{kārkhānyācā māl roj bājārāt jāto.}
\end{quote}

\begin{quote}
The factory's goods go to the market every day.
\end{quote}

Masculine.

\begin{cousinweb}
\textbf{\mr{माल}} (\emph{māl}) is a loan from Persian \emph{māl}, \textquotedblleft{}goods\textquotedblright{}, usually traced to Arabic \emph{māl}, \textquotedblleft{}property, wealth\textquotedblright{}.
\end{cousinweb}

\begin{grammarlens}[title={What you've built}]
One more word for describing what you lived and felt.
\end{grammarlens}

\subsection*{Your turn}
\begin{itemize}
\raggedright
  \item \emph{Say it:} \emph{māl}
  \item \emph{Say it:} \emph{māl}, once more
  \item \emph{Say it:} \emph{kārkhānyācā māl roj bājārāt jāto.} — The factory's goods go to the market every day.
  \item \emph{Recall:} say the Marathi for a factory, then the Marathi for production, then say \emph{māl} again
\end{itemize}

\begin{tcolorbox}[breakable,colback=teal!4,colframe=teal!35!black,title={Before you move on}]
What does \mr{माल} mean? (\textquotedblleft{}Goods; stock\textquotedblright{}.) Say it once more.
\end{tcolorbox}
\section[\texorpdfstring{\mr{पुरवठा} (puravṭhā)}{पुरवठा (puravṭhā)}]{\mr{पुरवठा} (puravṭhā) — a supply}

\label{lesson:MR-C512-puravtha}



\begin{quote}
Before the new one: say the Marathi for production, then the Marathi for goods.
\end{quote}

\subsection*{You'll want to know: \mr{पुरवठा}}
\textbf{\mr{पुरवठा}} — \emph{puravṭhā} — \textquotedblleft{}a supply\textquotedblright{}.

\begin{quote}
\mr{कारखान्याला} \mr{कच्च्या} \mr{मालाचा} \mr{पुरवठा} \mr{कमी} \mr{पडला}.
\end{quote}

\begin{quote}
\emph{kārkhānyālā kaccyā mālācā puravṭhā kamī paḍlā.}
\end{quote}

\begin{quote}
The factory ran short of raw materials.
\end{quote}

Masculine.

\mr{पुरवठा} \mr{कमी} \mr{पडणे}, to run short.

\begin{cousinweb}
\textbf{\mr{पुरवठा}} (\emph{puravṭhā}) is built on \mr{पुरवणे}, \textquotedblleft{}to provide\textquotedblright{}, the causative of \mr{पुरणे}, \textquotedblleft{}to suffice\textquotedblright{}, inherited from Sanskrit \emph{pūrayati}, \textquotedblleft{}fills\textquotedblright{} (CDIAL 8335).
\end{cousinweb}

\begin{grammarlens}[title={What you've built}]
One more word for describing what you lived and felt.
\end{grammarlens}

\subsection*{Your turn}
\begin{itemize}
\raggedright
  \item \emph{Say it:} \emph{puravṭhā}
  \item \emph{Say it:} \emph{puravṭhā}, once more
  \item \emph{Say it:} \emph{kārkhānyālā kaccyā mālācā puravṭhā kamī paḍlā.} — The factory ran short of raw materials.
  \item \emph{Recall:} say the Marathi for production, then the Marathi for goods, then say \emph{puravṭhā} again
\end{itemize}

\begin{tcolorbox}[breakable,colback=teal!4,colframe=teal!35!black,title={Before you move on}]
What does \mr{पुरवठा} mean? (\textquotedblleft{}A supply\textquotedblright{}.) Say it once more.
\end{tcolorbox}
\section[\texorpdfstring{\mr{खपणे} (khapṇe)}{खपणे (khapṇe)}]{\mr{खपणे} (khapṇe) — to sell well; to toil}

\label{lesson:MR-C512-khapne}



\begin{quote}
Before the new one: say the Marathi for goods, then the Marathi for a supply.
\end{quote}

\subsection*{You'll want to know: \mr{खपणे}}
\textbf{\mr{खपणे}} — \emph{khapṇe} — \textquotedblleft{}to sell well; to toil\textquotedblright{}.

\begin{quote}
\mr{सणांच्या} \mr{दिवसांत} \mr{सगळा} \mr{माल} \mr{खपला}.
\end{quote}

\begin{quote}
\emph{saṇāñcyā divsānt sagḷā māl khaplā.}
\end{quote}

\begin{quote}
In the festival days all the goods sold out.
\end{quote}

Past (he …): \mr{तो} \mr{माल} \mr{लवकर} \mr{खपला}, \emph{to māl lavkar khaplā}.

The goods are the subject: \mr{माल} \mr{खपतो}, the goods sell.

\begin{cousinweb}
\textbf{\mr{खपणे}} (\emph{khapṇe}) is a tadbhava: CDIAL lists Marathi \emph{khapṇẽ}, \textquotedblleft{}to be consumed, to work\textquotedblright{}, under Sanskrit \emph{kṣapyatē}, \textquotedblleft{}is destroyed\textquotedblright{} (CDIAL 3655); goods that sell are used up.
\end{cousinweb}

\begin{grammarlens}[title={What you've built}]
One more word for describing what you lived and felt.
\end{grammarlens}

\subsection*{Your turn}
\begin{itemize}
\raggedright
  \item \emph{Say it:} \emph{khapṇe}
  \item \emph{Say it:} \emph{khapṇe}, once more
  \item \emph{Say it:} \emph{saṇāñcyā divsānt sagḷā māl khaplā.} — In the festival days all the goods sold out.
  \item \emph{Recall:} say the Marathi for goods, then the Marathi for a supply, then say \emph{khapṇe} again
\end{itemize}

\begin{tcolorbox}[breakable,colback=teal!4,colframe=teal!35!black,title={Before you move on}]
What does \mr{खपणे} mean? (\textquotedblleft{}To sell well; to toil\textquotedblright{}.) Say it once more.
\end{tcolorbox}
